\documentclass[10pt,a4paper]{amsart}
\usepackage[margin=19mm]{geometry}
\usepackage{amsmath,amssymb,mathtools}
\usepackage[scr=rsfs,cal=euler]{mathalfa}
\usepackage{xfrac}
\usepackage[unicode,hidelinks,bookmarks=false]{hyperref}
\newcommand{\ie}{i.e.\ }
\newcommand{\cf}{cf.\ }
\newcommand{\resp}{respectively\ }
\newcommand{\Subsection}{\subsection}
\providecommand{\nobreakdash}{\nobreak\hbox{-}}
\setlength{\emergencystretch}{2em}
\multlinegap0pt
\usepackage{bm}
\usepackage{bbm}
\usepackage{dsfont}
\usepackage{xspace}
\usepackage{cancel}
\usepackage{mathtools}
\usepackage[shortlabels]{enumitem}
\usepackage{amsthm}
\makeatletter
\@ifundefined{theorem}{\newtheorem{theorem}{Theorem}}{}
\@ifundefined{definition}{\newtheorem{definition}[theorem]{Definition}}{}
\@ifundefined{proposition}{\newtheorem{proposition}[theorem]{Proposition}}{}
\makeatother
\makeatletter
\newcommand{\R}{\mathbb R}
\newcommand{\dd}{\,\mathrm{d}}
\expandafter\ifx\csname suggestedaddition\endcsname\relax
\newenvironment{suggestedaddition}
  {\par\begingroup\hypersetup{linkcolor=reviewgreen,citecolor=reviewgreen,urlcolor=reviewgreen}\color{reviewgreen}}
\fi
\expandafter\ifx\csname suggesteddeletion\endcsname\relax
\newenvironment{suggesteddeletion}
  {\par\begingroup\hypersetup{linkcolor=red,citecolor=red,urlcolor=red}\color{red}}
\fi
\makeatother
\title[A Harmonic-Map Proof of Mostow Rigidity]{A Harmonic-Map Proof of Mostow Rigidity}
\author{Guoyi Xu}
\date{Source: arXiv:2607.26677v1 (CC BY 4.0)}
\begin{document}
\maketitle

\noindent\emph{Source and licence.} A Harmonic-Map Proof of Mostow Rigidity. Version arXiv:2607.26677v1 (CC BY 4.0), distributed under the Creative Commons Attribution 4.0 International licence, as recorded in the arXiv licence metadata for this exact version and shown on the abstract page https://arxiv.org/abs/2607.26677v1. Full rights notice: licenses/arxiv-2607.26677-CC-BY-4.0.txt.\par\smallskip

\noindent\textit{Editorial scope.} Counted unit: the Proposition whose source label is \texttt{prop:elementary\_qs\_regular} in the article (source0.tex lines 1245--1271, source label \texttt{prop:elementary\_qs\_regular}), delivered together with the article's own complete original proof of it (source0.tex lines 1273--1619, one \texttt{proof} environment of 347 lines). No mathematics is written, repaired or re-derived: statement and proof are the article's own text line for line. Required context retained uncounted: def:quasisymmetry (lines 723--750). Every definition, display and label of those lines is retained unchanged. Omitted from this package and not redistributed: 1--722, 751--1244, 1272--1272, 1620--3351. Nothing inside a retained excerpt is omitted or altered: every retained body is delivered exactly as the article's lines are written, so no omission receipt of any kind accompanies this package. Citations: the retained excerpts carry no citation command at all, so no bibliography entry of the article is transcribed and no reference is redistributed. Reference adaptation: none was needed and none was made; every reference inside the delivered unit resolves inside it, so no unresolved reference or citation is produced. Printed slips kept literal: no printed word, symbol, hypothesis or number of the counted statement or of its proof was altered. Layout and class: the article's own class is \texttt{amsart} with packages including amssymb, amsmath, txfonts, mathrsfs, mathtools, titletoc, tikz, color, graphicx, float, enumitem, hyperref, cleveref; the wrapper is the lane's pinned \texttt{amsart} editorial wrapper at 10pt on A4, which restates the theorem-like environments the retained text needs and adds the article's own macro definitions that the retained text needs (\texttt{\textbackslash R}, \texttt{\textbackslash dd}), with extra wrapper packages (bm, bbm, dsfont, xspace, cancel, mathtools, enumitem). The delivered numbering is the unit's own. Assets: no retained span contains a figure environment, a graphics-inclusion command, a PDF-page inclusion, a table or a TikZ picture; no image file is copied into this package, the package declares no asset dependency and no image file is redistributed with it. Licence: version arXiv:2607.26677v1 of this article is distributed under the Creative Commons Attribution 4.0 International licence, as recorded in the arXiv licence metadata for this exact version and shown on the abstract page https://arxiv.org/abs/2607.26677v1. A full rights notice is delivered at licenses/arxiv-2607.26677-CC-BY-4.0.txt. Characters: every retained excerpt is pure ASCII, so the delivered engine, XeLaTeX with its default Latin Modern text font, reports no missing character.
\begin{definition}\label{def:quasisymmetry}
%\ReviewLabelChange{def:power-quasisym}{def:quasisymmetry}
Let \((X,d_X)\) and \((Y,d_Y)\) be metric spaces, and let
\begin{align}
\eta:[0,\infty)\longrightarrow[0,\infty)
\nonumber
\end{align}
be an increasing homeomorphism. A homeomorphism
\begin{align}
f:X\longrightarrow Y
\nonumber
\end{align}
is called \textbf{\(\eta\)-quasisymmetric} if
\begin{align}
\frac{d_Y\bigl(f(x),f(a)\bigr)}
     {d_Y\bigl(f(x),f(b)\bigr)}
\leq
\eta\!\left(
\frac{d_X(x,a)}{d_X(x,b)}
\right)
\label{eq:metric-quasisymmetry}
\end{align}
for all distinct \(x,a,b\in X\).

The map \(f\) is called \textbf{quasisymmetric} if it is
\(\eta\)-quasisymmetric for some increasing homeomorphism
\(\eta:[0,\infty)\to[0,\infty)\).
\end{definition}

\begin{proposition}\label{prop:elementary_qs_regular}
Let \(n\geq3\), let \(\Omega,\Omega'\subset\R^{n-1}\) be domains, and let
\begin{align}
f:\Omega\longrightarrow\Omega'\nonumber
\end{align}
be an \(\eta\)-quasisymmetric homeomorphism in the sense of
Definition~\ref{def:quasisymmetry}.  Then:
\begin{enumerate}[label=\textup{(\roman*)},leftmargin=3em]
\item
for every \(1\leq p<n-1\),
\begin{align}
f\in W^{1,p}_{\mathrm{loc}}(\Omega;\R^{n-1});\nonumber
\end{align}
\item
\(f\) is differentiable at Lebesgue-almost every point of
\(\Omega\), and its weak differential agrees almost everywhere with its
classical differential;
\item
at every point \(x\) where \(f\) is   differentiable,
\begin{align}
Df_x=0
\qquad\text{or}\qquad
Df_x\text{ is invertible}.
\label{eq:zero_or_invertible}
\end{align}
\end{enumerate}
\end{proposition}

\begin{proof}
\textbf{Step (1).} 
Fix bounded open sets
\begin{align}
V\Subset W\Subset\Omega.\nonumber
\end{align}
Choose \(r_0>0\) such that
\begin{align}
B(x,2r_0)\Subset W
\qquad
\text{for every }x\in V.\nonumber
\end{align}
Define a finite Borel measure on \(W\) by
\begin{align}
\nu(A):=\mathcal L^{n-1}(f(A)).\nonumber
\end{align}
Equivalently, \(\nu\) is the pushforward of Lebesgue measure on
\(f(W)\) by the continuous inverse \(f^{-1}\); hence it is indeed a
Borel measure.  Finiteness follows from \(\overline W\Subset\Omega\),
because \(f(\overline W)\) is compact.

Let \(x\in V\), let \(y\in\Omega\), and put
\begin{align}
r:=|x-y|<r_0.\nonumber
\end{align}
For every \(z\in\partial B(x,2r)\), quasisymmetry gives
\begin{align}
\frac{|f(x)-f(y)|}{|f(x)-f(z)|}
\leq
\eta(1/2).\nonumber
\end{align}
Since \(f\) is a homeomorphism and \(\overline{B(x,2r)}\Subset\Omega\),
\begin{align}
\partial f(B(x,2r))=f(\partial B(x,2r)).\nonumber
\end{align}
It follows that
\begin{align}
B\!\left(
 f(x),\frac{|f(x)-f(y)|}{\eta(1/2)}
\right)
\subset f(B(x,2r)).\nonumber
\end{align}
Indeed, the distance from \(f(x)\) to the boundary of the open set
\(f(B(x,2r))\) is at least the displayed radius.  Taking Euclidean
volumes, with \(\omega_{n-1}=\mathcal L^{n-1}(B(0,1))\), yields
\begin{align}
\omega_{n-1}
\left(
\frac{|f(x)-f(y)|}{\eta(1/2)}
\right)^{n-1}
\leq
\nu(B(x,2r)).\nonumber
\end{align}
Consequently,
\begin{align}
|f(x)-f(y)|
\leq
C_{n-1,\eta}|x-y|
\left(
\frac{\nu(B(x,2|x-y|))}{(2|x-y|)^{n-1}}
\right)^{1/(n-1)}.
\label{eq:qs_pointwise_density}
\end{align}

Define the truncated maximal density
\begin{align}
g(x)
:=
\left(
\sup_{0<s<2r_0}
\frac{\nu(B(x,s))}{s^{n-1}}
\right)^{1/(n-1)},
\qquad x\in V.
\label{eq:g_def}
\end{align}
The function \(g\) is measurable.  By continuity from below of \(\nu\),
the supremum may be taken over rational \(s\in(0,2r_0)\); for fixed
\(s\), the function \(x\mapsto\nu(B(x,s))\) is Borel, as follows by
approximating the ball indicator with continuous radial functions and
applying monotone convergence.
Then \eqref{eq:qs_pointwise_density} gives
\begin{align}
|f(x)-f(y)|
\leq
C_{n-1,\eta}|x-y|g(x)
\label{eq:hajlasz_one_point}
\end{align}
whenever \(x\in V\) and \(|x-y|<r_0\).

\textbf{Step (2).}
For the function \(g\) defined in \eqref{eq:g_def} and
\(\lambda>0\), put
\begin{align}
A_\lambda:=\{x\in V:g(x)>\lambda\}.\nonumber
\end{align}
For each \(x\in A_\lambda\), choose \(0<r_x<2r_0\) such that
\begin{align}
\nu(B(x,r_x))>\lambda^{n-1} r_x^{n-1}.\nonumber
\end{align}
The \(5r\)-covering lemma supplies a countable pairwise disjoint
subfamily
\begin{align}
B_i=B(x_i,r_i)\nonumber
\end{align}
such that
\begin{align}
A_\lambda\subset\bigcup_i5B_i.\nonumber
\end{align}
All \(B_i\) lie in \(W\).  Therefore
\begin{align}
\mathcal L^{n-1}(A_\lambda)
&\leq
5^{n-1}\omega_{n-1}\sum_i r_i^{n-1}\leq
\frac{5^{n-1}\omega_{n-1}}{\lambda^{n-1}}
\sum_i\nu(B_i)\leq
\frac{5^{n-1}\omega_{n-1}\nu(W)}{\lambda^{n-1}}.
\label{eq:g_weak_L_boundary}
\end{align}


Combining \eqref{eq:g_weak_L_boundary} with the trivial estimate
\(\mathcal L^{n-1}(A_\lambda)\leq\mathcal L^{n-1}(V)\), we have
\begin{align}
\mathcal L^{n-1}(A_\lambda)
\leq
\min\left\{
\mathcal L^{n-1}(V),
\frac{5^{n-1}\omega_{n-1}\nu(W)}
{\lambda^{n-1}}
\right\}.
\nonumber
\end{align}
For every \(p>0\), Tonelli's theorem gives
\begin{align}
\int_Vg^p\,\dd x
=
p\int_0^\infty
\lambda^{p-1}\mathcal L^{n-1}(A_\lambda)\,\dd\lambda.
\nonumber
\end{align}

Splitting the integral at \(\lambda=1\), for \(0<p<n-1\) we obtain
\begin{align}
\begin{aligned}
\int_Vg^p\,\dd x
&\leq
p\mathcal L^{n-1}(V)
\int_0^1\lambda^{p-1}\,\dd\lambda+
p\,5^{n-1}\omega_{n-1}\nu(W)
\int_1^\infty\lambda^{p-n}\,\dd\lambda\\
&=
\mathcal L^{n-1}(V)
+
\frac{
p\,5^{n-1}\omega_{n-1}\nu(W)
}{
n-1-p
}
<\infty.
\end{aligned}
\nonumber
\end{align}
Here the integral over \((0,1)\) is finite because \(p>0\), whereas
the integral over \((1,\infty)\) is finite because \(p<n-1\).
Thus
\begin{align}
g\in L^p(V)
\qquad
\text{for every }0<p<n-1.
\label{eq:g_Lp}
\end{align}

In particular, \(g(x)<\infty\) for almost every \(x\in V\), and
\eqref{eq:hajlasz_one_point} implies
\begin{align}
\varlimsup_{y\to x}
\frac{|f(y)-f(x)|}{|y-x|}
\leq C_{n-1,\eta}g(x)<\infty
\label{eq:finite_upper_lip}
\end{align}
for almost every \(x\in V\).

\textbf{Step (3).} We now prove directly that \eqref{eq:finite_upper_lip} implies  
differentiability almost everywhere.  For integers \(a,b\geq1\), let
\(E_{a,b}\subset V\) consist of all \(x\in V\) such that
\begin{align}
|f(x)|\leq a\nonumber
\end{align}
and
\begin{align}
|f(y)-f(x)|\leq a|y-x|
\label{eq:Eab_local_lip}
\end{align}
whenever \(y\in\Omega\) and \(0<|y-x|<1/b\).  These sets are measurable;
by continuity, the condition may be tested on a fixed countable dense
subset of \(\Omega\).  The set on which the upper pointwise Lipschitz
constant is finite is contained in
\begin{align}
\bigcup_{a,b\geq1}E_{a,b}.\nonumber
\end{align}

The restriction \(f|_{E_{a,b}}\) is globally Lipschitz.  Indeed, if
\(x,y\in E_{a,b}\) and \(|x-y|<1/b\), then
\eqref{eq:Eab_local_lip} applies; if \(|x-y|\geq1/b\), then
\begin{align}
|f(x)-f(y)|\leq2a\leq2ab|x-y|.\nonumber
\end{align}

For each nonempty
\(E_{a,b}\), let \(L_{a,b}:=\max\{a,2ab\}\).  For each coordinate function of
\(f|_{E_{a,b}}\), the McShane formula
\begin{align}
F^\alpha(z)
:=
\inf_{x\in E_{a,b}}
\bigl(f^\alpha(x)+L_{a,b}|z-x|\bigr)\nonumber
\end{align}
defines an \(L_{a,b}\)-Lipschitz extension to \(\R^{n-1}\).  Indeed, the
Lipschitz inequality for \(f^\alpha|_{E_{a,b}}\) shows that
\(F^\alpha=f^\alpha\) on \(E_{a,b}\), while the triangle inequality
shows that
\begin{align}
|F^\alpha(z)-F^\alpha(z')|\leq L_{a,b}|z-z'|.\nonumber
\end{align}
The vector-valued extension
\begin{align}
F=(F^1,\dots,F^{n-1})\nonumber
\end{align}
is \(\sqrt{n-1}\,L_{a,b}\)-Lipschitz, since every coordinate is
\(L_{a,b}\)-Lipschitz.  Rademacher's theorem therefore makes \(F\)
  differentiable almost everywhere.

Let \(x\in E_{a,b}\) be both a density point of \(E_{a,b}\) and a
  differentiability point of \(F\).  We claim that \(f\) itself is
differentiable at \(x\), with derivative \(DF_x\).  Given any sequence
\(y_k\to x\), the density of \(E_{a,b}\) at \(x\) allows us to choose
\(z_k\in E_{a,b}\) such that
\begin{align}
|z_k-y_k|=o(|y_k-x|).
\label{eq:zk_approx_yk}
\end{align}
Indeed, otherwise some \(\varepsilon>0\) and a subsequence would
satisfy
\begin{align}
B(y_k,\varepsilon|y_k-x|)\cap E_{a,b}=\varnothing,\nonumber
\end{align}
which would contradict that \(x\) is a density point.  

For large \(k\),
\eqref{eq:Eab_local_lip} at \(z_k\) and
\eqref{eq:zk_approx_yk} give
\begin{align}
|f(y_k)-f(z_k)|
\leq
a|y_k-z_k|
=
o(|y_k-x|).\nonumber
\end{align}
Since \(F=f\) on \(E_{a,b}\) and \(F\) is differentiable at \(x\),
\begin{align}
f(z_k)-f(x)
=
DF_x(z_k-x)+o(|z_k-x|).\nonumber
\end{align}
Together with \eqref{eq:zk_approx_yk}, this yields
\begin{align}
f(y_k)-f(x)
=
DF_x(y_k-x)+o(|y_k-x|).\nonumber
\end{align}
Thus \(f\) is   differentiable at \(x\).  The Lebesgue density
theorem, Rademacher's theorem, and the countable union over \(a,b\) show
that \(f\) is   differentiable almost everywhere in \(V\).
Since \(V\Subset\Omega\) was arbitrary, the same holds in \(\Omega\).

\textbf{Step (4).} Fix \(1\leq p<n-1\) and \(K\Subset V\).  For a coordinate vector \(e_i\)
and sufficiently small nonzero \(h\), define
\begin{align}
D_{i,h}f(x)
:=
\frac{f(x+he_i)-f(x)}{h},
\qquad x\in K.\nonumber
\end{align}
By \eqref{eq:hajlasz_one_point},
\begin{align}
|D_{i,h}f(x)|\leq C_{n-1,\eta}g(x)\nonumber
\end{align}
for almost every \(x\in K\).  At almost every differentiability point,
\begin{align}
D_{i,h}f(x)\longrightarrow Df_x(e_i)
\qquad(h\to0),\nonumber
\end{align}
and the same bound holds for \(|Df_x(e_i)|\).  Since \(g\in L^p(K)\) by \eqref{eq:g_Lp},
dominated convergence gives
\begin{align}
D_{i,h}f
\longrightarrow
Df(e_i)
\qquad
\text{in }L^p(K;\R^{n-1}).
\label{eq:difference_quotient_Lp}
\end{align}
For \(\psi\in C_c^\infty(K)\) and sufficiently small \(h\), a change of
variables gives
\begin{align}
\int_KD_{i,h}f(x)\psi(x)\,\dd x
=
\int_\Omega
f(x)
\frac{\psi(x-he_i)-\psi(x)}{h}
\,\dd x.\nonumber
\end{align}
Letting \(h\to0\) and using \eqref{eq:difference_quotient_Lp}, we obtain
\begin{align}
\int_KDf_x(e_i)\psi(x)\,\dd x
=
-\int_Kf(x)\partial_i\psi(x)\,\dd x.\nonumber
\end{align}
Hence \(Df(e_i)\) is the \(i\)-th weak derivative.  This proves
\begin{align}
f\in W^{1,p}_{\mathrm{loc}}(\Omega;\R^{n-1})\nonumber
\end{align}
and shows that the weak and classical differentials agree almost
everywhere.

\textbf{Step (5).} Zero or invertible.
Let \(x\in\Omega\) be a differentiability point, and let
\(v,w\in\mathbb S^{n-2}\).  For sufficiently small \(t\neq0\), the points
\(x+tv\) and \(x+tw\) are equally distant from \(x\).  Quasisymmetry gives
\begin{align}
|f(x+tv)-f(x)|
\leq
\eta(1)|f(x+tw)-f(x)|.\nonumber
\end{align}
Divide by \(|t|\) and let \(t\to0\):
\begin{align}
|Df_x(v)|
\leq
\eta(1)|Df_x(w)|.
\label{eq:linear_distortion_at_derivative}
\end{align}
If \(Df_x(w)=0\) for one unit vector \(w\), then
\eqref{eq:linear_distortion_at_derivative} implies \(Df_x(v)=0\) for every
unit vector \(v\), so \(Df_x=0\).  Otherwise \(Df_x\) has trivial kernel;
because its domain and range both have dimension \(n-1\), it is invertible.
This proves \eqref{eq:zero_or_invertible}.
\end{proof}

\end{document}
